\documentclass[11pt]{article}
\usepackage{amsmath,amssymb,amsthm}
\usepackage{hyperref}

\newtheorem{theorem}{Theorem}
\newtheorem{proposition}[theorem]{Proposition}
\newtheorem{lemma}[theorem]{Lemma}
\newtheorem{conjecture}[theorem]{Conjecture}
\newtheorem{definition}[theorem]{Definition}
\newtheorem{remark}[theorem]{Remark}
\newtheorem{question}[theorem]{Question}

\newcommand{\N}{\mathbb{N}}
\newcommand{\Z}{\mathbb{Z}}
\newcommand{\Q}{\mathbb{Q}}
\newcommand{\CPP}{\operatorname{CPP}}
\newcommand{\GK}{\operatorname{GK}}
\newcommand{\QSym}{\operatorname{QSym}}
\newcommand{\Sym}{\operatorname{Sym}}
\newcommand{\Des}{\operatorname{Des}}
\newcommand{\charge}{\operatorname{charge}}
\newcommand{\SSYT}{\operatorname{SSYT}}
\newcommand{\wt}{\operatorname{wt}}

% q-series notation
\newcommand{\qpoch}[2]{(#1;#2)_\infty}
\newcommand{\qpochn}[3]{(#1;#2)_{#3}}
\newcommand{\qbin}[2]{\genfrac{[}{]}{0pt}{}{#1}{#2}}

\title{Literature connections: word statistics, charge,\\
and $L$-positivity for cylindric partitions}
\author{Notes for research}
\date{\today}

\begin{document}
\maketitle

\begin{abstract}
We collect four threads from the literature that connect
to the $L$-positivity of the cylindric partition frequency
generating function $\widetilde{\GK}_c \in \QSym$
established in our companion note (QUASI\_SYMMETRIC.tex).
The goal is to identify the right combinatorial framework
for understanding this positivity and its relationship
to Warnaar's Conjecture~2.7.
\end{abstract}

%% ========================================================================
\section{The charge statistic (Lascoux--Sch\"utzenberger)}
\label{sec:charge}
%% ========================================================================

Lascoux and Sch\"utzenberger~\cite{LS78,LS79} introduced the
\emph{charge} statistic on words and semistandard Young tableaux.
For a tableau $T \in \SSYT(\lambda, \mu)$ (shape~$\lambda$,
content~$\mu$), the charge $\charge(T) \in \Z_{\geq 0}$ is
computed from the \emph{reading word} $w(T)$ as follows:

\begin{enumerate}
  \item Read the entries of~$T$ from right to left, bottom to top,
    to obtain the word $w(T) = w_1 w_2 \cdots w_n$.
  \item Decompose $w(T)$ into ``standard subwords'' by
    successively extracting the leftmost occurrence of
    $1, 2, \ldots, \mu_1$, then repeating.
  \item The charge of each standard subword is
    $\sum_{i \in \Des(\sigma)} (n - i)$, where $\sigma$ is
    the subword read as a permutation and $\Des$ denotes
    its descent set.
  \item $\charge(T)$ is the sum of the charges of all subwords.
\end{enumerate}

The main theorem is that Kostka--Foulkes polynomials satisfy
\[
  K_{\lambda,\mu}(q) = \sum_{T \in \SSYT(\lambda,\mu)} q^{\charge(T)}.
\]

\begin{question}
Is there an analog of the charge statistic for cylindric
partitions? Specifically, given a CPP
$\pi = (\nu^{(1)}, \nu^{(2)}, \nu^{(3)})$, can one define
a reading word $w(\pi)$ and a charge-like statistic such that
the $Q_{n,c}(q)$ polynomials (or the $L$-coefficients of
$\widetilde{\GK}_c$) are given by a charge-weighted sum?
\end{question}

\begin{remark}
We have previously verified computationally that $Q_{n,c}(q)$
does \emph{not} decompose as a nonneg sum of Kostka--Foulkes
polynomials $K_{\lambda,\mu}(q)$. However, the QSym framework
may require a \emph{different} charge-like statistic adapted to
compositions rather than partitions.
\end{remark}

%% ========================================================================
\section{The plactic monoid and crystal bases}
\label{sec:plactic}
%% ========================================================================

The \emph{plactic monoid} (Lascoux--Sch\"utzenberger~\cite{LS81})
is the quotient of the free monoid on positive integers by the
Knuth relations:
\begin{align*}
  xzy &\equiv zxy \quad \text{if } x \leq y < z, \\
  yxz &\equiv yzx \quad \text{if } x < y \leq z.
\end{align*}
Two words are \emph{plactically equivalent} if and only if they
have the same RSK insertion tableau. The plactic monoid provides
a combinatorial model for the representation ring of $GL_n$:
\[
  s_\lambda \cdot s_\mu = \sum_\nu c^\nu_{\lambda\mu}\, s_\nu,
\]
where the Littlewood--Richardson coefficient $c^\nu_{\lambda\mu}$
counts plactic classes of words with certain content and
RSK shape conditions.

In the language of Kashiwara's crystal bases, the plactic monoid
is the monoid of the crystal graph $B(\infty)$ for type~$A$.
The crystal operators $\tilde{e}_i, \tilde{f}_i$ act on words
and preserve the RSK recording tableau while modifying the
insertion tableau.

\begin{question}
The descent set of a word is a plactic invariant (it depends only
on the RSK insertion tableau via the descent set of the
corresponding standard tableau). Since $\widetilde{\GK}_c$ is
$L$-positive and $L_\alpha$ encodes descent sets, can the
$L$-coefficients $b_\alpha(c)$ be interpreted as counting
plactic classes of ``cylindric words'' with descent
composition~$\alpha$?
\end{question}

\begin{remark}
For type $A_2$ (rank~3), the crystal graph has a natural
connection to the three partitions $(\nu^{(1)}, \nu^{(2)},
\nu^{(3)})$ in a CPP via the weight lattice of~$\mathfrak{sl}_3$.
The cylindric condition (cyclic interlacing) may correspond to
a periodicity condition on crystal paths, connecting to the
\emph{affine} crystal $\widehat{B}(\Lambda)$ for
$\widehat{\mathfrak{sl}}_3$.
\end{remark}

%% ========================================================================
\section{Keys and Demazure atoms (Lascoux)}
\label{sec:keys}
%% ========================================================================

Lascoux~\cite{Las90} introduced \emph{keys} (also called
\emph{Demazure characters} or \emph{key polynomials})
$\kappa_\alpha$ indexed by compositions $\alpha$, and their
``atomic'' decomposition: the \emph{Demazure atoms}
$\mathfrak{A}_\alpha$ satisfy
\[
  \kappa_\alpha = \sum_{\beta \leq \alpha} \mathfrak{A}_\beta,
\]
where $\leq$ is the Bruhat-like order on compositions.

The Demazure atoms form a basis of the polynomial ring
$\Z[x_1, \ldots, x_n]$ (in finitely many variables) that is
closely related to the fundamental quasi-symmetric functions:
under the limit $n \to \infty$, the atoms stabilize to
the $L_\alpha$.

More precisely, Mason~\cite{Mas09} showed that
the \emph{quasi-key polynomials} (also called
\emph{Demazure atoms}) map to fundamental quasi-symmetric
functions under the forgetful map from polynomials to
quasi-symmetric functions.

\begin{question}
Can $\widetilde{\GK}_c$ be lifted from $\QSym$ to the polynomial
ring as a positive sum of Demazure atoms? If so, the extra
structure (dependence on a specific number of variables) might
encode the max-entry information that is lost in the QSym
framework.
\end{question}

%% ========================================================================
\section{$P$-partitions and labeled posets (Stanley, Gessel)}
\label{sec:P-partitions}
%% ========================================================================

Stanley's theory of $P$-partitions~\cite{Sta72} provides the
foundational connection between labeled posets and quasi-symmetric
functions. Gessel~\cite{Ges84} showed that the generating function
for $P$-partitions of a labeled poset $(P, \omega)$ is always
a nonneg sum of fundamental quasi-symmetric functions:
\[
  \Gamma_{(P,\omega)}(\mathbf{z})
  = \sum_{\substack{f: P \to \Z_{>0} \\ f \text{ is a
    $(P,\omega)$-partition}}} \prod_{p \in P} z_{f(p)}
  = \sum_\alpha c_\alpha\, L_\alpha,
  \qquad c_\alpha \in \Z_{\geq 0}.
\]

A \emph{$(P,\omega)$-partition} is an order-preserving map
$f: P \to \Z_{>0}$ that satisfies strict inequalities
$f(p) < f(q)$ whenever $p <_P q$ and $\omega(p) > \omega(q)$
(i.e., at ``descents'' of the labeling).

\begin{question}\label{q:poset}
Can cylindric partitions be realized as $P$-partitions for
a suitable labeled poset? Concretely:
\begin{enumerate}
  \item Define a poset $P_c$ whose elements are the positions
    $(i, j)$ (partition~$i$, row~$j$) of a CPP with profile~$c$.
  \item The interlacing conditions $\nu^{(i)}_j \geq
    \nu^{(i')}_{j'}$ define order relations on~$P_c$.
  \item A labeling $\omega: P_c \to \{1, 2, \ldots\}$ determines
    which inequalities are strict vs.\ weak.
  \item A CPP is then exactly a $(P_c, \omega)$-partition,
    and $L$-positivity of $\widetilde{\GK}_c$ follows from
    Stanley--Gessel theory.
\end{enumerate}
\end{question}

The key subtlety is the \emph{cyclic} nature of the interlacing:
$\nu^{(3)}_j \geq \nu^{(1)}_{j+c_0}$ closes the cycle
$1 \to 2 \to 3 \to 1$. A standard (acyclic) poset cannot directly
model this. However, by unrolling the cycle (working with a
finite number of values), the resulting poset may be acyclic,
and the $P$-partition framework may apply.

\begin{remark}
If Question~\ref{q:poset} has a positive answer, then the
$L$-positivity of $\widetilde{\GK}_c$ would follow from
general theory, and the $L$-coefficients $b_\alpha$ would
count \emph{linear extensions} of~$P_c$ with descent
set~$D(\alpha)$. This would provide a concrete
combinatorial interpretation of the $b_\alpha$.
\end{remark}

%% ========================================================================
\section{Summary and connections}
\label{sec:summary}
%% ========================================================================

The four threads above suggest the following picture:

\begin{center}
\begin{tabular}{l|l|l}
  \textbf{Framework} & \textbf{Key object} & \textbf{Connection to CPPs} \\
  \hline
  Charge (LS) & $q^{\charge(T)}$ on SSYT
    & Refines $K_{\lambda,\mu}(q)$; \\
    & & need cylindric analog \\
  Plactic (LS) & Crystal graph $B(\lambda)$
    & Affine crystals for $\widehat{\mathfrak{sl}}_3$ \\
    & & model cylindric tableaux \\
  Keys (Lascoux) & Demazure atoms $\mathfrak{A}_\alpha$
    & Lift from QSym to polynomials; \\
    & & may encode max-entry \\
  $P$-partitions (SG) & Labeled poset $(P, \omega)$
    & CPP interlacing $\approx$ poset; \\
    & & cyclic closure is the obstacle
\end{tabular}
\end{center}

The most promising direction for proving $L$-positivity appears
to be the $P$-partition approach (Section~\ref{sec:P-partitions}),
as it would give $L$-positivity directly from general theory.
The charge/crystal approaches may be more relevant for the
original $Q_{n,c}(q)$ positivity problem, where the max-entry
variable must be incorporated.

\begin{thebibliography}{99}
\bibitem{Ges84}
I.~Gessel,
\emph{Multipartite $P$-partitions and inner products of skew
Schur functions},
Contemp.\ Math.\ \textbf{34} (1984), 289--317.

\bibitem{Las90}
A.~Lascoux,
\emph{Polyn\^omes de Schubert: une approche historique},
S\'eminaire Lotharingien de Combinatoire \textbf{B25a} (1990).

\bibitem{LS78}
A.~Lascoux and M.-P.~Sch\"utzenberger,
\emph{Sur une conjecture de H.\,O.~Foulkes},
C.\,R.\ Acad.\ Sci.\ Paris S\'er.~A \textbf{286} (1978), 323--324.

\bibitem{LS79}
A.~Lascoux and M.-P.~Sch\"utzenberger,
\emph{Le mono\"ide plaxique},
Quad.\ Ricerca Sci.\ \textbf{109} (1981), 129--156.

\bibitem{LS81}
A.~Lascoux and M.-P.~Sch\"utzenberger,
\emph{Croissance des polyn\^omes de Foulkes--Green},
C.\,R.\ Acad.\ Sci.\ Paris S\'er.~A \textbf{288} (1979), 95--98.

\bibitem{Mas09}
S.~Mason,
\emph{An explicit construction of type~A Demazure atoms},
J.~Algebraic Combin.\ \textbf{29} (2009), no.~3, 295--313.

\bibitem{Sta72}
R.\,P.~Stanley,
\emph{Ordered structures and partitions},
Mem.\ Amer.\ Math.\ Soc.\ \textbf{119} (1972).
\end{thebibliography}

\end{document}
